\documentclass{article}
\usepackage{amsmath}
\usepackage{hyperref}
\usepackage{graphicx}
\usepackage{minted}
\usepackage[htt]{hyphenat}
\usepackage[UTF8]{ctex}

\hypersetup{
  colorlinks=true,
  allcolors=blue,
  pdfstartview=Fit,
  breaklinks=true
}

\newcommand{\figph}[2][0.8\textwidth]{%
  \IfFileExists{figures/#2}%
    {\includegraphics[width=#1]{figures/#2}}%
    {\fbox{\parbox[c][4cm][c]{0.75\textwidth}{\centering\ttfamily [图片占位符]\quad\detokenize{figures/#2}}}}%
}

\sloppy

\begin{document}
  \title{系统开发工具基础 \\ 第四周实验报告}
  \author{乔海轩—25090012035}
  \date{2026年9月14日}
  \maketitle

  \tableofcontents
  \newpage

  \section{总览}
    \begin{table}[h]
      \centering
      \begin{tabular}{|l|l|c|}
        \hline
        题目编号 & 描述 & 完成情况 \\
        \hline
        13 & 建立可执行的本地质量门禁 & 已完成 \\
        14 & 让 Make 只重建真正受影响的产物 & 已完成 \\
        15 & 把本地 API 数据转换为可读报告 & 已完成 \\
        16 & 修复并交付一个陌生的小型工具仓库 & 已完成 \\
        \hline
      \end{tabular}
    \end{table}
    Github 仓库地址：\href{https://github.com/Ptilotsis-o/system-dev-tools/}{https://github.com/Ptilotsis-o/system-dev-tools/}
    \newline Commit记录如图\ref{fig:commit1}和\ref{fig:commit2}所示。
    \begin{figure}[h]
      \centering
      \figph{commit1.png}
      \caption{Commit记录截图}
      \label{fig:commit1}
    \end{figure}
    \begin{figure}[h]
      \centering
      \figph{commit2.png}
      \caption{Commit记录截图}
      \label{fig:commit2}
    \end{figure}

  \newpage
  \section{第 13 题 建立可执行的本地质量门禁}
    \subsection*{题目要求}
      \begin{itemize}
        \item 复制 q10 为 q13；在 pyproject.toml 中加入 ruff 和 pytest 的最小配置。
        \item 补充一个正常姓名测试和一个空白姓名测试，每个测试包含有效断言。
        \item 运行 ruff format、ruff check 和 pytest，修复全部问题，不得全局忽略规则。
        \item 编写 check.sh，依次执行 ruff format --check、ruff check 和 pytest；运行并保证返回 0。
      \end{itemize}
    \subsection*{操作步骤与命令}
      \begin{minted}{bash}
        cp -r ../week3/q10 q13
        cd q13
        vim pyproject.toml
        vim tests/test_cli.py
        ruff format .
        ruff check .
        pytest -v
        vim check.sh
        chmod +x check.sh
        ./check.sh
        echo $?
      \end{minted}
    \subsection*{结果验证}
      \begin{itemize}
        \item pyproject.toml 中加入 ruff 最小配置（target-version、line-length、lint 的 select 列表，ignore 为空）与 pytest 最小配置（testpaths 为 tests）；
        \item \texttt{test\_normal\_name} 与 \texttt{test\_blank\_name\_exits} 各自包含有效断言，均通过；
        \item \texttt{ruff format}、\texttt{ruff check} 无任何报错，未使用全局忽略规则；
        \item check.sh 依次执行 \texttt{ruff format --check}、\texttt{ruff check} 与 \texttt{pytest}，运行结束打印 All checks passed，\texttt{echo \$?} 输出 0。
      \end{itemize}
      \begin{figure}[h]
        \centering
        \figph{ruff.png}
        \caption{ruff format , ruff check 与 pytest 输出截图}
        \label{fig:ruff}
      \end{figure}
      \begin{figure}[h]
        \centering
        \figph{checksh.png}
        \caption{check.sh 运行及退出码截图}
        \label{fig:checksh}
      \end{figure}

  \newpage
  \section{第 14 题 让 Make 只重建真正受影响的产物}
    \subsection*{题目要求}
      \begin{itemize}
        \item 创建给定的 data.csv、stats.py、build\_report.py 和 report.md。
        \item 编写 Makefile，包含 all、stats.txt、report.txt 和 clean；准确列出数据文件、脚本和中间产物依赖，并把 all、clean 声明为 .PHONY。
        \item 验证首次 make 生成两个产物；第二次无改动时不执行配方。
        \item touch data.csv 后再次 make，确认 stats.txt 和 report.txt 重新生成；clean 只删除生成物。
      \end{itemize}
    \subsection*{操作步骤与命令}
      \begin{minted}{bash}
        vim data.csv stats.py build_report.py report.md
        vim Makefile
        make
        make
        touch data.csv
        make
        make clean
        ls
      \end{minted}
    \subsection*{结果验证}
      \begin{itemize}
        \item 依赖关系准确：\texttt{stats.txt} 依赖 \texttt{data.csv} 与 \texttt{stats.py}，\texttt{report.txt} 依赖 \texttt{report.md} 与 \texttt{stats.txt}；all、clean 已声明为 \texttt{.PHONY}；
        \item 首次 \texttt{make} 由 data.csv 计算得到 stats.txt（内容为 10），再生成 report.txt；
        \item 第二次 \texttt{make} 未执行任何配方，只提示 \texttt{'stats.txt' is up to date}；
        \item \texttt{touch data.csv} 后再次 \texttt{make}，stats.txt 与 report.txt 均重新生成，时间戳更新；
        \item \texttt{make clean} 只删除 stats.txt 与 report.txt，源文件与 Makefile 保留。
      \end{itemize}
      \begin{figure}[h]
        \centering
        \figph{make.png}
        \caption{首次 make 生成两个产物截图和第二次 make 不执行配方截图}
        \label{fig:make}
      \end{figure}
      \begin{figure}[h]
        \centering
        \figph{makeagain.png}
        \caption{touch 后重新生成截图}
        \label{fig:makeagain}
      \end{figure}

  \newpage
  \section{第 15 题 把本地 API 数据转换为可读报告}
    \subsection*{题目要求}
      \begin{itemize}
        \item 把给定 JSON 保存为 packages.json，并在 q15 目录运行 python -m http.server 8000。
        \item 使用 curl -fsS 获取 http://127.0.0.1:8000/packages.json。
        \item 使用 jq 筛选 status 为 active 且 downloads 不少于 100 的记录，并按 downloads 降序、name 升序排列。
        \item 编写 api\_report.sh，把结果生成 summary.md，包含标题和 name、version、downloads 三列表格。
      \end{itemize}
    \subsection*{操作步骤与命令}
      \begin{minted}{bash}
        vim packages.json
        python -m http.server 8000 &
        curl -fsS http://127.0.0.1:8000/packages.json
        vim api_report.sh
        chmod +x api_report.sh
        ./api_report.sh
        cat summary.md
      \end{minted}
    \subsection*{结果验证}
      \begin{itemize}
        \item packages.json 保存正确，\texttt{python -m http.server 8000} 在 q15 目录提供该文件；
        \item \texttt{curl -fsS} 成功取回 JSON 内容，未依赖浏览器或手工复制；
        \item jq 筛选出 alpha、gamma、delta 三条记录（beta 为 inactive、epsilon 下载量不足 100 被排除），排序为 downloads 降序、同值按 name 升序：delta（450）、gamma（450）、alpha（120）；
        \item summary.md 包含标题 \texttt{\# Active Packages} 与 name、version、downloads 三列表格，共三行数据，与 api\_report.sh 的输出一致。
      \end{itemize}
      \begin{figure}[h]
        \centering
        \figph{server.png}
        \caption{http.server截图}
        \label{fig:server}
      \end{figure}
      \begin{figure}[h]
        \centering
        \figph{crul.png}
        \caption{crul捕获截图}
        \label{fig:curl}
      \end{figure}
      \begin{figure}[h]
        \centering
        \figph{summary.png}
        \caption{summary.md 内容截图}
        \label{fig:summary}
      \end{figure}

  \newpage
  \section{第 16 题 修复并交付一个陌生的小型工具仓库}
    \subsection*{题目要求}
      \begin{itemize}
        \item 复制 q13 为 q16 并初始化/继续使用 Git。
        \item 把问候语实现临时改成始终返回字面量 Hello, name!，运行 check.sh 确认测试失败；随后定位并修复。
        \item 编写最小 Makefile，包含 check、build 和 clean；check 调用 check.sh、build 生成 wheel。
        \item 运行 make check 和 make build，计算 wheel 的 SHA-256，并把最终改动提交为一个内容聚焦的 Git 提交。
      \end{itemize}
    \subsection*{操作步骤与命令}
      \begin{minted}{bash}
        cp -r q13 q16
        cd q16
        git init
        git add -A
        git commit -m "chore: import greetlab from q13"
        vim src/greetlab/cli.py # 临时改动
        ./check.sh
        vim src/greetlab/cli.py # 修复
        ./check.sh
        vim Makefile
        make check
        make build
        sha256sum dist/greetlab_25090012035-0.1.0-py3-none-any.whl
        git add -A
        git commit -m "···"
      \end{minted}
    \subsection*{结果验证}
      \begin{itemize}
        \item q16 由 q13 复制而来并纳入 Git 管理；
        \item 临时改动后 check.sh 返回非 0，\texttt{test\_normal\_name} 断言失败，实际输出为 \texttt{Hello, name!} 而非 \texttt{Hello, Alice!}；
        \item 修复后 \texttt{make check} 返回 0，ruff format --check、ruff check 与 pytest 全部通过；
        \item \texttt{make build} 生成 \allowbreak\texttt{greetlab\_25090012035-0.1.0-py3-none-any.whl}，并计算出该 wheel 的 SHA-256；
        \item 最终改动以单个内容聚焦的提交入库（\texttt{fix(cli): restore blank-name validation}）。
      \end{itemize}
      \begin{figure}[h]
        \centering
        \figph{fail.png}
        \caption{临时改动后 check.sh 测试失败截图}
        \label{fig:fail}
      \end{figure}
      \begin{figure}[h]
        \centering
        \figph{check.png}
        \caption{make check 输出截图}
        \label{fig:makecheck}
      \end{figure}
      \begin{figure}[h]
        \centering
        \figph{sha.png}
        \caption{wheel 的 SHA-256 计算结果截图}
        \label{fig:sha}
      \end{figure}
      \begin{figure}[h]
        \centering
        \figph{q16commit.png}
        \caption{最终 Git 提交记录截图}
        \label{fig:q16commit}
      \end{figure}

  \newpage
  \section{课后作业}
    \subsection*{q1: 为字符串工具包建立质量门禁}
    目标：创建一个字符串工具包，实现字符串反转和回文判断功能；配置 ruff 和 pytest；编写 check.sh，依次运行 ruff format --check、ruff check 和 pytest，确保代码格式、静态检查和单元测试全部通过且返回 0。
      \begin{minted}{bash}
        vim pyproject.toml
        pip install -e .
        vim check.sh
        chmod +x check.sh
        ./check.sh
        echo $?
      \end{minted}
      \subsubsection*{结果验证}
        \begin{itemize}
          \item \texttt{reverse()} 实现字符串反转（\texttt{"abc" → "cba"}），\texttt{is\_palindrome()} 实现回文判断（忽略大小写和空格）；
          \item ruff 配置了 line-length = 88，lint select 为 E、F、W，无全局忽略规则；pytest 配置 testpaths 为 tests；
          \item check.sh 返回 0，ruff format --check、ruff check 与 pytest 全部通过，三个测试均通过。
        \end{itemize}
        \begin{figure}[h]
          \centering
          \figph{check1.png}
          \caption{q1 check.sh 测试通过截图}
          \label{fig:homework-check}
        \end{figure}
    \subsection*{q2: 用 Makefile 构建 Markdown 文档流水线}
    目标：创建多个 Markdown 章节文件，用 Python 脚本合并并生成目录；编写 Makefile 管理文件依赖，实现首次构建生成最终文档，无改动时不重复执行，修改源文件后触发级联重建，并能通过 clean 目标只删除生成物。
      \begin{minted}{bash}
        vim Makefile
        make
        make
        touch chapter1.md
        make
        make clean
      \end{minted}
      \subsubsection*{结果验证}
        \begin{itemize}
          \item \texttt{combine.py} 合并所有 chapter*.md 生成 book.md，\texttt{toc.py} 提取标题生成目录插入后得到 final.md；
          \item \texttt{make} 第一次构建生成 book.md 和 final.md，第二次构建提示 \texttt{'final.md' is up to date}，不重复执行；
          \item \texttt{touch chapter1.md} 后，\texttt{make} 级联重建 book.md 和 final.md，时间戳更新；
          \item \texttt{make clean} 仅删除 book.md 和 final.md，源文件与脚本保留。
        \end{itemize}
        \begin{figure}[h]
          \centering
          \figph{make1.png}
          \caption{q2 第一和第二次 make 输出结果截图}
          \label{fig:homework-make1}
        \end{figure}
        \begin{figure}[h]
          \centering
          \figph{make2.png}
          \caption{q2 \texttt{touch chapter1.md} 后 make 输出结果截图}
          \label{fig:homework-make2}
        \end{figure}
    \subsection*{q3: 用 jq 生成 CSV 报告}
    目标：创建商品库存 JSON 文件，启动本地 HTTP 服务；编写脚本用 curl 获取数据，通过 jq 筛选出库存小于 10 的商品，按价格降序排列，输出为包含表头的 CSV 文件。
      \begin{minted}{bash}
        vim inventory.json
        vim report.sh
        chmod +x report.sh
        python -m http.server 8080 &
        ./report.sh
        cat inventory_low.csv
        kill %1
      \end{minted}
      \subsubsection*{结果验证}
        \begin{itemize}
          \item \texttt{report.sh} 通过 curl 获取 JSON，用 jq 筛选 \texttt{stock < 10} 的记录，按 price 降序排列，输出带表头的 CSV；
          \item 筛选结果：desk（stock=3, price=120.0）、chair（stock=5, price=45.0）、apple（stock=8, price=5.0）、notebook（stock=7, price=2.0），共 4 条记录，按价格降序排列正确。
        \end{itemize}
        \begin{figure}[h]
          \centering
          \figph{jq.png}
          \caption{q3 report.sh 输出结果截图}
          \label{fig:homework-jq}
        \end{figure}
    \subsection*{q4: 修复并打包一个小型计算器 CLI}
    目标：创建一个计算器命令行工具，故意在减法操作中植入加法缺陷；通过测试发现失败并修复；编写 Makefile 管理 check 和 build；构建 wheel 并计算 SHA-256；初始化 Git 并提交最终修复。
      \begin{minted}{bash}
        mkdir -p src/calc tests
        vim src/calc/cli.py          # 减法故意写为加法
        vim tests/test_cli.py
        vim pyproject.toml
        pip install -e .
        pytest -v                     # test_sub 失败
        vim src/calc/cli.py           # 修复为 args.a - args.b
        pytest -v
        vim Makefile check.sh
        make check
        make build
        sha256sum dist/calc_25090012035-0.1.0-py3-none-any.whl
        # git init && git add -A && git commit -m "fix(calc): restore subtraction"
      \end{minted}
      \subsubsection*{结果验证}
        \begin{itemize}
          \item 首次运行 pytest，\texttt{test\_sub} 输出 \texttt{3.0}（实际应为 5+2=7.0），断言失败；
          \item 修复减法为 \texttt{args.a - args.b} 后，\texttt{test\_add} 和 \texttt{test\_sub} 全部通过；
          \item \texttt{make check} 运行 check.sh（ruff format --check、ruff check、pytest）全部通过，返回 0；
          \item \texttt{make build} 生成 \texttt{calc\_25090012035-0.1.0-py3-none-any.whl}，SHA-256 计算成功；
          \item Git 提交信息为 \texttt{fix(calc): restore subtraction}，内容聚焦。
        \end{itemize}
        \begin{figure}[h]
          \centering
          \figph{wrong.png}
          \caption{q4 测试失败}
          \label{fig:homework-wrong}
        \end{figure}
        \begin{figure}[h]
          \centering
          \figph{calc.png}
          \caption{q4 测试修复}
          \label{fig:homework-calc}
        \end{figure}
        \begin{figure}[h]
          \centering
          \figph{hash.png}
          \caption{q4 构建 wheel SHA-256 计算结果}
          \label{fig:homework-hash}
        \end{figure}
    \subsection*{q5: 禁止 print 的质量门禁}
    目标：创建一个日志工具模块，其中使用 print 输出；配置 ruff 启用 T201 规则禁止 print；将 print 替换为 logging；编写测试并集成到 check.sh，确保 ruff 检查和 pytest 全部通过。
      \begin{minted}{bash}
        vim logger.py                  # 先使用 print 输出
        ruff check .                   # 触发 T201 报错
        vim logger.py                  # 将 print 替换为 logging.info
        vim tests/test_logger.py
        vim pyproject.toml
        vim check.sh
        chmod +x check.sh
        ./check.sh
        echo $?
      \end{minted}
      \subsubsection*{结果验证}
        \begin{itemize}
          \item 初始 \texttt{print(msg)} 被 ruff T201 规则检测并报错；
          \item 替换为 \texttt{logging.info(f"LOG: \{msg\}")} 后，ruff 检查通过；
          \item \texttt{test\_log} 使用 caplog fixture 验证日志输出，pytest 通过；
          \item check.sh 返回 0，ruff 检查与 pytest 全部通过。
        \end{itemize}
        \begin{figure}[h]
          \centering
          \figph{ruff1.png}
          \caption{q5 ruff 检查失败}
          \label{fig:homework-ruff}
        \end{figure}
        \begin{figure}[h]
          \centering
          \figph{print.png}
          \caption{q5 check.sh 测试通过截图}
          \label{fig:homework-print}
        \end{figure}
    \subsection*{q6: Makefile 管理图片缩略图生成}
    目标：创建多个 txt 文件，编写缩略图生成脚本；使用 Makefile 的模式规则为每个 txt 文件生成对应的 thumb.txt；实现首次全部生成，新增文件后只生成新增的缩略图，并通过 clean 删除所有缩略图。
      \begin{minted}{bash}
        mkdir images
        vim images/a.txt images/b.txt images/c.txt
        vim thumbnail.py
        vim Makefile
        make                    # 首次生成所有缩略图
        make                    # 无改动不执行
        vim images/d.txt
        make                    # 只生成 d.thumb.txt
        make clean              # 删除所有 .thumb.txt
      \end{minted}
      \subsubsection*{结果验证}
        \begin{itemize}
          \item \texttt{thumbnail.py} 读取源 txt 文件，生成以 \texttt{THUMB: } 为前缀的缩略文件；
          \item Makefile 使用模式规则 \texttt{\%.thumb.txt: \%.txt}，准确列出依赖；
          \item 首次 \texttt{make} 生成 a.thumb.txt、b.thumb.txt、c.thumb.txt；
          \item 第二次 \texttt{make} 不执行任何配方；新增 d.txt 后只生成 d.thumb.txt；
          \item \texttt{make clean} 仅删除 images/*.thumb.txt，源文件保留。
        \end{itemize}
        \begin{figure}[h]
          \centering
          \figph{thumb1.png}
          \caption{q6 第一次缩略图生成结果截图}
          \label{fig:homework-thumb1}
        \end{figure}
        \begin{figure}[h]
          \centering
          \figph{thumb2.png}
          \caption{q6 第二次缩略图生成结果截图}
          \label{fig:homework-thumb2}
        \end{figure}
    \subsection*{q7: jq 分组统计并输出 Markdown 列表}
    目标：创建销售记录 JSON 文件，使用 jq 按地区分组统计总销售额，按总销售额降序排列，输出为 Markdown 列表格式并保存为 sales\_summary.md。
      \begin{minted}{bash}
        vim sales.json
        vim summarize.sh
        chmod +x summarize.sh
        ./summarize.sh
        cat sales_summary.md
      \end{minted}
      \subsubsection*{结果验证}
        \begin{itemize}
          \item \texttt{summarize.sh} 使用 jq 的 \texttt{group\_by(.region)} 按地区分组，\texttt{map(.amount) | add} 求和，\texttt{sort\_by(-.total)} 降序排列；
          \item 输出结果：North（1500）、South（1200）、East（200），格式为 Markdown 无序列表 \texttt{- region: total}；
          \item sales\_summary.md 内容与预期一致。
        \end{itemize}
        \begin{figure}[h]
          \centering
          \figph{summary1.png}
          \caption{q7 sales\_summary.md 内容截图}
          \label{fig:homework-sales}
        \end{figure}
    \subsection*{q8: 增加 mypy 类型检查到质量门禁}
    目标：创建一个小型数学库，故意写错类型注解；配置 mypy 和 pytest；编写 check.sh 依次运行 mypy、ruff check 和 pytest；修复类型错误后确保所有检查通过。
      \begin{minted}{bash}
        vim mathlib.py               # 故意写错类型注解
        vim tests/test_mathlib.py
        vim pyproject.toml
        mypy .                       # 触发类型错误
        vim mathlib.py               # 修复注解
        mypy .                       # 通过
        vim check.sh
        chmod +x check.sh
        ./check.sh
        echo $?
      \end{minted}
      \subsubsection*{结果验证}
        \begin{itemize}
          \item \texttt{mathlib.py} 中 \texttt{add(a: int, b: int) -> int} 故意写错（如 b 类型错误），mypy 检测到类型不匹配并报错；
          \item 修复后 mypy 通过，无任何类型错误；
          \item check.sh 依次执行 mypy、ruff check、pytest，全部通过并返回 0。
        \end{itemize}
        \begin{figure}[h]
          \centering
          \figph{mypy.png}
          \caption{q8 check.sh 测试通过截图}
          \label{fig:homework-mypy}
        \end{figure}
    \subsection*{q9: Makefile 构建 Python wheel 并运行测试}
    目标：创建一个简单 Python 包，编写 Makefile，定义 all 目标生成 wheel，test 目标运行测试，clean 目标清理生成物；实现首次构建生成 wheel，无改动时不重复构建，并确保测试通过。
      \begin{minted}{bash}
        mkdir -p src/demo tests
        vim src/demo/core.py
        vim tests/test_core.py
        vim pyproject.toml
        vim Makefile
        make                 # 首次构建生成 wheel
        make                 # 无改动不重复构建
        make test            # 运行测试
        make clean           # 清理生成物
      \end{minted}
      \subsubsection*{结果验证}
        \begin{itemize}
          \item \texttt{src/demo/core.py} 实现 \texttt{greet(name) → "Hello, \{name\}!"}，pytest 测试通过；
          \item Makefile 中 all 依赖 \texttt{dist/demo-0.1.0-py3-none-any.whl}，依赖 src/demo/core.py 和 pyproject.toml；
          \item 首次 \texttt{make} 构建生成 wheel；第二次 \texttt{make} 提示 up to date，不重复执行；
          \item \texttt{make test} 先构建 wheel 再运行 pytest，全部通过；
          \item \texttt{make clean} 删除 dist、build、*.egg-info 和 \_\_pycache\_\_。
        \end{itemize}
        \begin{figure}[h]
          \centering
          \figph{wheel.png}
          \caption{q9 测试通过截图}
          \label{fig:homework-wheel}
        \end{figure}
    \subsection*{q10: jq 修改 JSON 并生成差异报告}
    目标：创建配置文件 JSON，使用 jq 将所有版本号加 1 生成新文件；用 diff 生成差异报告文件；编写脚本封装该流程，并确保脚本最终返回 0。
      \begin{minted}{bash}
        vim config.json
        vim update_versions.sh
        chmod +x update_versions.sh
        ./update_versions.sh
        cat config.diff
        echo $?
      \end{minted}
      \subsubsection*{结果验证}
        \begin{itemize}
          \item \texttt{update\_versions.sh} 使用 \texttt{jq 'map(.version += 1)'} 将每个服务的版本号加 1，输出到 config\_new.json；
          \item config.diff 正确展示版本号变更（service-a: 1→2, service-b: 2→3, service-c: 1→2）以及 JSON 格式化差异；
          \item 脚本最终返回 0，流程完整。
        \end{itemize}
        \begin{figure}[h]
          \centering
          \figph{diff.png}
          \caption{q10 config.diff 差异报告截图}
          \label{fig:homework-diff}
        \end{figure}
    \subsection*{课后作业 commit 记录}
      \begin{figure}[h]
        \centering
        \figph{homework_commit.png}
        \caption{课后作业 commit 记录}
        \label{fig:homework_commit}
      \end{figure}

  \section{解题感悟}
    本次实验报告围绕第四周的工程化实践展开，涵盖了代码质量门禁、Makefile 增量构建、jq 数据处理以及 Git 综合交付。通过完成课后作业，我不仅掌握了 ruff、pytest、mypy 等工具的配置与集成，还理解了依赖管理、自动化检查和数据流水线的设计思路。在综合实例中，我进一步锻炼了从修复缺陷到构建 wheel、计算哈希并提交版本的全流程能力。总的来说，这次实践让我对开发工具链的协同工作有了更具体的认识，也为今后规范化交付项目积累了经验。

\end{document}